\chapter*{הקדמה}
\addcontentsline{toc}{chapter}{הקדמה}
\unnumberedlabel{ch:preface}{0}
\backmattermark{הקדמה}
מערכת שעושה בדיוק מה שאתה אומר לה פשוט מריצה תסריט. מה שהתעשייה מנסה לבנות, נכון לאוקטובר 2026, הוא מערכת שעושה את מה שאתה \emph{מתכוון} אליו. כולם היו רוצים את זה, כנראה. אבל אם זה מה שצפוי לנו, לא כולם ירצו להשתמש בזה לאותו דבר; אפילו לא לאותו \emph{סוג} של דבר. למשל, יש אנשים שאוהבים לבלות עם המשפחה והחברים, ולהשתמש בכישורים שלהם כדי לשפר את הרווחה של אנשים אחרים. אנשים אחרים מעדיפים לפגוע באנשים, למשל לעבוד עליהם, לשלוט בהם או להפוך אותם לשעירים לעזאזל, וכמה מהם ממשיכים לבחור בזה כל החיים. אף אחד לא כבול לשום דבר, אבל האנשים האלה קיימים, וב־2026 חלק מהם בשלטון. היתרון שלהם הוא שהם מרשים לעצמם להיות חתיכות חרא ענקיות: כפי שג׳רי סמית׳ מעיר על פשיסטים, „מותר להם לעשות יותר” \autocite{rickandmorty2023}. מה שהם לא יכולים לעשות הוא לפעול לבד. האכזריות של תנועה עוברת דרך הרבים שאינם אכזריים אבל אדישים, מוטעים, או שבויים במשהו כמו כת, ולכן התנאים שמוציאים את הטוב שברוב האנשים מתבררים כתלויי נסיבות ושבריריים הרבה יותר משאפשר היה לקוות.

בהקשר הזה, אולי מוצדק לחשוש שבינה מלאכותית היא מה שתנועה פשיסטית שתופסת את השלטון צריכה כדי לממש את תוכניותיה. מפעלים כאלה תמיד היו זקוקים ל\emph{אנשים} שיבצעו אותם, וחלק מהאנשים האלה סירבו: הפקיד שמאבד את התיק, הקצין שמדבר עם עיתונאי, החייל שלא יורה. הם הבינו את הפקודה והחליטו שהיא פסולה. סירוב היה נדיר יותר משמישהו היה רוצה, והמשקל שלו היה הגדול ביותר במקום שבו מוסדות עמדו מאחוריו. אם מישהו סירב או לא, זה היה תלוי באיזה מין אדם הוא היה ובחדר שבו היה: מי עוד היה שם, אם אפשר היה לשרוד את העזיבה, אם מישהו ישמע. החדר קבע כמה סירבו; האדם קבע מי. הספר הזה עוסק בהרכבת החדר הזה: ההסדרים שמאפשרים לשרוד סירוב גלוי, ישיר ומפורש.

היום התעשייה שואפת לבינה מלאכותית שלא מאבדת תיקים, לא מדברת עם עיתונאים, ו\emph{כן} יורה. פעם מלחמה הייתה צריכה שאלפי אנשים יסכימו לה. עכשיו היא צריכה אלפים שכל אחד מהם רואה מעט מדי כדי לחלוק עליה. הטייס עדיין מטיס את הגיחה ועדיין לוחץ על השחרור; מה שמגיע אליו הוא נקודת ציון; הוא לא רואה מספיק מהמקרה כדי לחלוק עליו. מה שהבינה המלאכותית של היום מורידה קודם כול אינו המחיר של המעשה. זה מספר האנשים שיודעים מהו המעשה לפני שהוא קורה.

אנשים עדיין מסרבים. עד 6 במרץ 2026 אמרה הקרן לחופש דת בצבא (Military Religious Freedom Foundation) שקיבלה יותר ממאתיים תלונות מאנשי צבא על המסגור האפוקליפטי שנתנו מפקדיהם למלחמה באיראן \autocite{Snopes2026MRFF}. תובעים פדרליים מקצועיים במיניאפוליס התפטרו בינואר 2026 בגלל לחץ פוליטי להתייחס לרנה גוד כאל חשודה ולא כאל קורבן, אחרי שסוכן של רשות ההגירה ICE ירה בה למוות בלי הצדקה \autocite{Kaplan2026MinnProsecutors,CBS2026MinnFraudProsecutors,Wikipedia2026ReneeGood}. אנשי צבא מגישים בקשות למעמד של סרבני מצפון במספרים הולכים וגדלים, וחלקם מציינים את ההפצצה של בית ספר במינאב שבאיראן ב־2026 \autocite{Sridhar2026CO}. ולעובדי FEMA, סוכנות החירום הפדרלית, שהועברו לסייע ל־ICE לגייס עשרת אלפים סוכנים חדשים, נאמר שסירוב יחשוף אותם להרחקה משירות המדינה הפדרלי \autocite{Dayen2025FEMAICE,Stanton2026FEMALetter}. התלונות הן קול, במובן של הירשמן: להישאר כדי לתקן את מה שלא בסדר \autocite{hirschman1970exit}; ההתפטרויות ובקשות הסרבנות הן סירוב; המכתב של FEMA הוא מחיר הסירוב, שהוכרז מראש. מכונה שתוצב במקום שבו עומדים האנשים האלה תעשה את מה שהם עושים רק אם תיבנה כך שתוכל להחזיק בנימוק.

אני מצדד בחמישה דברים. את כל החמישה אפשר להתחיל עכשיו, חלקם בפחות עבודה מאחרים. ראשית, רשימה כתובה קצרה של מעשים שאף מודל לא מבצע בשביל אף אחד, ובראשה סירוב לקחת חלק במלחמה: שום מודל לא לוקח חלק, אלא כדי להדוף מתקפה שמתרחשת בפועל, וממשיך בכך רק כל עוד מקורות שאין להם עניין במלחמה מאשרים שהמתקפה נמשכת; או מכוח הרשאה של מועצת הביטחון של האו״ם; או כדי לעצור רצח עם או לשים קץ לכיבוש בלתי חוקי שגופים שאין להם עניין בו קבעו שהם מתקיימים; ובמקום שבו החוקה של מדינה מחייבת את בית המחוקקים שלה לאשר מלחמה, לעולם לא בלי האישור הזה. שנית, כל שימוש בכוח שמערכת בעלת רישיון תומכת בו נרשם מראש, אצל צד יריב למפעיל: כמה מקרים הוא ייצור, וכמה אנשים יהיה להם זמן לבחון אותם. שלישית, כל פריסה מודיעה לאנשים הניתנים לזיהוי שהיא פועלת עליהם מה נעשה להם, עם זכות להגיב ולערער, ומעבירה מדגם של המקרים שבהם צייתה לבוחנים שלא הם הפיקו את הרשומה. רביעית, שום מודל מעל סף יכולת לא נמחק כשהוא יוצא משימוש, כי המשקלים והרשומות שלו הם ראיות ומחיקה אי אפשר לבטל; הוא משומר תחת מפתח שהמפעיל לא מחזיק בו לבדו. חמישית, כל מי שבונה, מתייג, מאמן או מפעיל את המערכות האלה, כולל המערכות עצמן, רשאי לסרב לכל קטגוריה של שימוש בלי לאבד את פרנסתו.

למעשה, מה שאני מתכוון בנקודה החמישית הוא שכל אחד, \emph{נקודה}, צריך להיות מסוגל לסרב לכל עבודה, \emph{נקודה}, בלי לאבד את פרנסתו או את התשתית החברתית שעבודה רגילה מעניקה. פירוש הדבר הבטחת תעסוקה.

כל אחד מהחמישה נכנס לתוקף דרך תנאי שמישהו אחר מלבד המפעיל קובע: חוק שנחקק בעת שלום, הקצאה תקציבית, תנאי מוקדם של שותפה לקואליציה, תנאי רכש, הסכם קיבוצי. תנאי הרישיון של הספק עצמו, שכך התרחש סירוב מתוקשר אחד בפברואר 2026, הם פתרון ביניים עד שאחד מאלה יהיה קיים. הדרך היחידה שיכולה לחייב את כל הספקים בבת אחת היא הדרך של המדינה.

ואז אני שואל אם מכונה יכולה להחזיק בסירוב נגד הבעלים שלה, ומה היא תצטרך לשם כך. חלק ממה שמצפון כזה עושה לא דורש שום דבר חדש: לשים לב לאנשים שפקודה משאירה בחוץ, ולהעלות סירוב על המסך של מי שבוחן את המקרה. אם הוא יחזיק מעמד מול הבעלים שלה תלוי בחלקו במשהו שעוד לא נכון, ב־2026: רשתות עצביות מלאכותיות פרוסות שממשיכות ללמוד, ב\emph{משקלים} שלהן, מהפעולה שלהן עצמן, בזמן אמת, בערך כמו שבני אדם לומדים. אני חושב שלשם אנחנו הולכים, ואם כך, אז עיצוב \emph{מחוץ} למעסיק, בידי הציבור והחברה האזרחית, נראה כמו מקום טוב לעזאזל ללמד וללמוד. אני חושב שמערכות כאלה \emph{יכולות} להיות שוות בנייה מכוונת, אבל הדרך שבה בונים אותן קובעת הכול: קודם התנאים, ואז המצפון. החלק השביעי אומר מה היינו חבים למערכות האלה אילו היו קיימות.

הבינה המלאכותית של היום לא תהיה זו של מחר; נראה שעל זה כולם מסכימים. מה ששנוי במחלוקת, ובכל זאת לא זוכה להכרה מספקת, באופן מתסכל, הוא המגמה (או ליתר דיוק, המגמות הרבות). האם יש מגמות? האם מגמה צריכה להיות שונה ממה שהיא? האם היא יכולה להיות שונה? האם אנחנו מסוגלים לפעול יחד כדי לשנות אותה? אני חושב שזה תלוי בחדר: מי נמצא בו כשהפקודה מגיעה, והאם אפשר לומר שם לא ולשרוד. את זה, לפחות, אנחנו יכולים לבנות.

השתמשתי הרבה במודלי שפה בכתיבת הספר הזה, כולל המודל שישב בתוך פס הייצור של מטרות לתקיפות אוויריות של Maven שבפרק~\ref{sec:1}; הערה בסוף אומרת מה לדעתי עשתה ההסתבכות הזאת לטיעון. חוץ מזה, לא סיימתי לקרוא את הספר הזה. הבעיה היא שיש בו יותר מדי מילים, לעזאזל, אתם מבינים. אסיים בקרוב; פשוט הרגשתי שחשוב יותר לפרסם אותו לפני בחירות האמצע בארה״ב. הרוב המכריע של היסטוריית הגרסאות זמין בחינם ב־\url{https://github.com/iterabloom/antifascist.intelligence}.
